\section{Design}
\label{sec:method}

\textsc{Prefetch} records, for every function, the working set of pages that the
function accessed during its most recent invocations. On a cold start it restores
the snapshot and loads the recorded working set in one sequential read before the
function resumes, instead of serving each page on demand.

\paragraph{Recording.} During an invocation, the hypervisor tracks first-touch
page faults. The recorded working set $W_f$ of function $f$ is the union of the
pages touched in the last $k$ invocations:
\begin{equation}
  W_f = \bigcup_{i=n-k+1}^{n} P_f^{(i)},
  \label{eq:working-set}
\end{equation}
where $P_f^{(i)}$ is the set of pages touched in invocation $i$. We use $k=3$.

\paragraph{Prefetching.} At restore time, \textsc{Prefetch} reads the pages of
$W_f$ from the snapshot file in offset order and installs them before the vCPU
starts. Pages outside $W_f$ are still served lazily, so a mispredicted working
set costs latency but never correctness.

Figure~\ref{fig:pipeline} summarizes the restore path.

\begin{figure}[t]
  \centering
  \setlength{\unitlength}{1mm}
  \begin{picture}(120,22)
    \put(0,6){\framebox(30,10){snapshot}}
    \put(30,11){\vector(1,0){12}}
    \put(42,6){\framebox(32,10){read $W_f$}}
    \put(74,11){\vector(1,0){12}}
    \put(86,6){\framebox(32,10){resume vCPU}}
  \end{picture}
  \caption{Restore path in \textsc{Prefetch}: the recorded working set is read
    sequentially before the function resumes.}
  \label{fig:pipeline}
\end{figure}
